% ============================================================================
\appendix
\section{Proofs}\label{app:proofs}
% ============================================================================

\begin{proof}[Proof of Proposition~\ref{prop:exact}]
Fix $v$ and $k$ and let $R^\star=\{i:z_{ivk}=1\}$. Constraint \eqref{eq:route} gives $\sum_R y_{Rvk}=1$ and $\sum_{R\ni i}y_{Rvk}=z_{ivk}$. For $i\in R^\star$ the second equality equals the first, so every $R$ with $y_{Rvk}>0$ contains $i$; for $i\notin R^\star$ the sum is zero, so no such $R$ contains $i$. Hence $y_{Rvk}>0$ implies $R=R^\star$, and $y_{R^\star vk}=1$. The routing cost is $\mathrm{TSP}(R^\star)$, vehicle loads are at most $Q$ by \eqref{eq:load}, and each store is on at most one vehicle. Every feasible lookahead plan (stores per vehicle, loads, inventory flows) has exactly one representation of this form, which proves exactness.
\end{proof}

\begin{proof}[Proof of Theorem~\ref{thm:fifo}]
(i) Setting $z=0$, $y_{\emptyset vk}=1$, $q=0$, $\zeta=0$ and serving demand from available stock in any order is feasible: day-0 stock satisfies the shelf constraint by assumption and stock only decreases afterwards. All cost terms (travel, holding of carried units, waste, lost sales) are non-negative, so the objective is bounded below by zero; with finitely many assignments and a bounded linear programme for each, an optimum exists.

(ii) Take an optimal recourse solution and view inventory as fungible flows of units per cohort. \emph{Greedy service.} If on day $k$ demand is lost while $\varepsilon>0$ units are left on the shelf, sell $\varepsilon$ of them. Those units originally had one of four fates: sold on a later day (that later sale is now lost; the cost change is $-p_u\varepsilon+p_u\varepsilon-h\varepsilon\cdot(\text{nights saved})\le0$), expired or disposed (saving $p_s\varepsilon$ and holding), or carried beyond the horizon (saving holding). The cost does not increase. \emph{FIFO.} If on day $k$ a younger cohort $r'$ is sold while $\varepsilon$ units of an older cohort $r<r'$ are left, sell $\varepsilon$ older units instead and keep $\varepsilon$ younger ones. The kept younger units replicate the original path of the older units: they are sellable on every day the older units were (their remaining life is longer); where the older units expired at the end of day $k+r-1$, the younger units are disposed on that day at the same cost $p_s$ and with the same number of holding nights; disposal and carrying beyond the horizon are replicated identically. The number of units carried each night is unchanged, so the cost is unchanged. Each exchange reduces $\sum_k\sum_r r\,w_{irk}$ and the feasible set is a polytope, so a finite sequence of exchanges yields an optimal solution with FIFO issuing and greedy service.

(iii) $q$ enters the recourse LP only through the right-hand side of $S_{imk}=q_k$, so $V_{\mathrm{LP}}$ is the value function of an LP in its right-hand side, which is convex and piecewise linear. The physical FIFO trajectory with $\zeta=0$ is feasible for the recourse LP and has cost $C_{\mathrm{FIFO}}$; hence $V_{\mathrm{LP}}\le C_{\mathrm{FIFO}}$. By (ii) there is an optimal solution with FIFO issuing and greedy service; if it uses no disposal it coincides with the physical trajectory and the two values are equal.
\end{proof}

\begin{proof}[Proof of Proposition~\ref{prop:marginal}]
(a) By LP duality, the optimal dual of the constraint $S_{imk}=q_k$ is a subgradient of the value function with respect to $q_k$. Adding $\varepsilon$ to $q_k$ and disposing of these units immediately is feasible and costs $p_s\varepsilon$ (holding $h\varepsilon$ through $S_{imk}$ plus $(p_s-h)\varepsilon$ through $\zeta$), so $V_{\mathrm{LP}}(q+\varepsilon e_k)\le V_{\mathrm{LP}}(q)+p_s\varepsilon$ and $\pi_k\le p_s$. Removing $\varepsilon\le q_k$ units turns at most $\varepsilon$ sales into lost sales and cannot increase holding or waste, so $V_{\mathrm{LP}}(q-\varepsilon e_k)\le V_{\mathrm{LP}}(q)+p_u\varepsilon$ and $\pi_k\ge -p_u$.

(b) $C_{\mathrm{FIFO}}$ is a composition of finitely many $\min$, $\max$ and affine maps of $q$, hence piecewise linear and Lipschitz. Compare the FIFO trajectories with deliveries $q$ and $q+\varepsilon e_k$. On day $k$ the extra units are the freshest and are sold last, so sales of older stock are unchanged. On later days the extra units are older than subsequent deliveries and may displace them, but stock that was older than the extra units is never affected. Therefore, relative to the original trajectory, lost sales can only fall, waste and carried stock can only rise, and the total extra stock $E_\ell$ carried on night $\ell$ satisfies $0\le E_\ell\le\varepsilon$. Each extra unit leaves the system by preventing a lost sale ($-p_u$), by expiring on some day $\tau\le H-1$ ($+p_s$ after at most $\tau-k\le H-1-k$ nights), or by remaining at the end ($H-k$ nights). The cost change per unit is therefore at least $-p_u$ and at most $\max\{p_s+(H-1-k)h,\,(H-k)h\}=p_s+(H-1-k)h$ when $h\le p_s$. Dividing by $\varepsilon$ gives \eqref{eq:marg} for both one-sided derivatives.
\end{proof}

\begin{proof}[Proof of Proposition~\ref{prop:dfl}]
(a) With the visit pattern fixed and no capacity binding, the coverage map gives $q_k=\sum_{j\in W_k}\tilde d_j-U_k(\tilde d)$, where $U_k$ is the planned usable stock that serves the window. A unit increase of $\tilde d_j$, $j\in W_k$, increases $q_k$ by one unless the planned stock already covers day $j$, in which case the target of day $j$ does not affect any delivery; days outside every window affect no delivery. The chain rule then gives \eqref{eq:stationary}, and setting it to zero is the first-order condition of $\min_q \mathbb E[C_{\mathrm{FIFO}}(q;D)\mid\xi]$ over the deliveries of the fixed pattern.

(b) With the trigger pattern fixed and no capacity binding, a triggered delivery is $q_{it}=\tilde S_{it}-U_{it}$, where the usable stock $U_{it}$ depends on earlier deliveries and demands but not on $\tilde S_{it}$. Hence $\partial q_{it}/\partial\tilde S_{it}=1$, and $\tilde S_{it}$ affects the episode only through $q_{it}$; later deliveries respond to it through their usable stock, with their own levels held fixed. The chain rule gives \eqref{eq:stationaryul}, with $\mathrm dC^{T}/\mathrm dq_{it}$ the total derivative along the trajectory. If $\delta_\theta$ can take any value for each input, stationarity of $\mathbb E[C^{T}]$ in $\theta$ implies that \eqref{eq:stationaryul} vanishes for every input, which is the first-order condition of minimising $\mathbb E[C^{T}]$ over the order-up-to levels of the triggered store-days given their inputs.

(c) With $m=1$ nothing is carried: $W_k=\{k\}$ and $q_k=\tilde d_k$ in (a), $U_{it}=0$ and $q_{it}=\tilde S_{it}$ in (b), and the cost of a day, $p_u(D-q)^++p_s(q-D)^+$, does not depend on other days. Thus $\mathbb E[\partial C/\partial q\mid\xi]=-p_u\mathbb P(D>q\mid\xi)+p_s\mathbb P(D<q\mid\xi)$, which vanishes when $F_{D\mid\xi}(q)=p_u/(p_u+p_s)$ (for continuous $F$; for discrete demand, at the smallest $q$ with $F\ge p_u/(p_u+p_s)$).
\end{proof}

% ============================================================================
\section{First versions of the rule tuning and of ScenPO-SS, and the rule as published}\label{app:variants}
% ============================================================================

Table~\ref{tab:variants} compares, on the earlier test episodes of every configuration, the first versions of three policies with the versions reported in the main text. UL$_\lambda$ tuned on the fine grid is up to 0.6\% cheaper than when tuned on the coarse grid, significantly in six configurations. About half of the advantage of the first version of UL-DFL, trained from the coarse rule, therefore came from the resolution of the grid: against the rule tuned on the fine grid it was significantly cheaper in six configurations and more expensive with a four-day shelf life. Retrained from the fine rule, UL-DFL is 0.8--1.2\% cheaper than it in nine configurations. ScenPO-SS with the scenario set of \citet{jia2025scenario} is more expensive than the first implementation in twelve configurations (significantly in nine) and cheaper with a two-day shelf life. The first implementation sampled ten scenarios with equal probability from the twelve bands of the full quantile grid; three of these bands lie above the 0.9 quantile, where only 9\% of the probability mass lies, so its scenarios over-represented high demand. It served 92.0--97.1\% of demand, against 87.2--94.7\% for the published scenario set, and with expensive routing it would have been the cheapest policy (18,197 per episode against 18,630 for LA-Quantile2). On the fresh base-case test episodes the two versions of ScenPO-SS differ by 0.8\%, which is not significant.

The two rows on the published rule compare UL$_\lambda$ with the deliver-up-to-level rule exactly as published by \citet{crama2018stochastic}: one target service level for both the trigger and the delivery quantity, the total stock counted, and $\lambda\in\{1,m-1\}$, with the service level tuned on the same validation episodes over the fine grid (the paper's 90\% is in the grid). The published rule is significantly more expensive than our UL$_\lambda$ in every configuration, by 0.6--3.4\% (1.6\% on the fresh base-case test episodes), and UL-DFL is 0.9--3.8\% cheaper than it. Tuning the two service levels separately and netting out the stock about to expire therefore make the rule a stronger comparator, not a weaker one.

\begin{table}[H]
\centering\footnotesize
\caption{First versus final versions, and the rule as published, on the earlier test episodes: paired cost difference of the first-named policy relative to the second, in percent. Bold: 95\% paired bootstrap interval excludes zero.}\label{tab:variants}
\resizebox{\textwidth}{!}{\begin{tabular}{lrrrrrrrrrrrrr}
\toprule
Comparison (\%) & Base & Clustered & Mixed & $m=2$ & $m=4$ & $p_u=10$ & $p_u=40$ & $\rho=1.10$ & $\rho=1.60$ & High var. & $c^{\mathrm{r}}=3$ & 30\% LIFO & Real data \\
\midrule
UL$_\lambda$, fine vs.\ coarse grid & \textbf{-0.5} & +0.1 & +0.1 & -0.3 & \textbf{-0.6} & +0.0 & -0.6 & \textbf{-0.5} & \textbf{-0.6} & -0.4 & -0.2 & \textbf{-0.6} & \textbf{-0.5} \\
UL-DFL vs.\ UL$_\lambda$, both coarse (first version) & \textbf{-0.8} & \textbf{-1.2} & \textbf{-1.0} & \textbf{-0.8} & -0.1 & \textbf{-1.0} & \textbf{-1.1} & \textbf{-0.9} & \textbf{-0.5} & -0.4 & \textbf{-0.8} & \textbf{-0.7} & \textbf{-0.5} \\
UL-DFL from coarse rule vs.\ UL$_\lambda$ fine & -0.3 & \textbf{-1.2} & \textbf{-1.1} & \textbf{-0.5} & \textbf{+0.4} & \textbf{-1.0} & \textbf{-0.4} & -0.3 & +0.1 & -0.0 & \textbf{-0.6} & -0.1 & +0.0 \\
UL-DFL vs.\ UL$_\lambda$, both fine (reported) & \textbf{-0.8} & \textbf{-1.2} & \textbf{-1.0} & \textbf{-1.1} & -0.3 & \textbf{-1.0} & \textbf{-0.8} & \textbf{-0.8} & -0.3 & -0.0 & \textbf{-0.8} & \textbf{-0.8} & \textbf{+0.6} \\
ScenPO-SS, Jia et al.\ vs.\ band sampling & \textbf{+2.2} & +1.7 & \textbf{+3.5} & \textbf{-1.8} & \textbf{+7.6} & \textbf{+5.0} & \textbf{+5.0} & +0.5 & \textbf{+10.6} & \textbf{+3.5} & \textbf{+6.3} & +0.3 & \textbf{+3.6} \\
UL$_\lambda$ as published vs.\ UL$_\lambda$ (reported) & \textbf{+1.7} & \textbf{+0.8} & \textbf{+0.6} & \textbf{+0.6} & \textbf{+2.4} & \textbf{+1.5} & \textbf{+3.4} & \textbf{+1.3} & \textbf{+1.4} & \textbf{+0.9} & \textbf{+1.5} & \textbf{+1.9} & \textbf{+2.9} \\
UL-DFL vs.\ UL$_\lambda$ as published & \textbf{-2.4} & \textbf{-2.0} & \textbf{-1.6} & \textbf{-1.7} & \textbf{-2.6} & \textbf{-2.4} & \textbf{-3.8} & \textbf{-2.0} & \textbf{-1.6} & \textbf{-0.9} & \textbf{-2.2} & \textbf{-2.6} & \textbf{-2.1} \\
\midrule
Fill rate (\%), ScenPO-SS, Jia et al. & 92.4 & 92.7 & 92.0 & 92.1 & 92.4 & 88.0 & 94.7 & 91.9 & 91.5 & 87.2 & 91.1 & 92.3 & 89.2 \\
Fill rate (\%), ScenPO-SS, band sampling & 95.6 & 96.1 & 95.7 & 95.0 & 95.7 & 92.0 & 97.1 & 93.8 & 96.0 & 92.0 & 94.2 & 95.5 & 93.1 \\
\bottomrule
\end{tabular}
}
\end{table}
